% !TEX program = uplatex
\documentclass[uplatex,dvipdfmx,12pt,a4paper]{jsarticle}
\usepackage[margin=1in]{geometry}
\usepackage{amsmath,amssymb,mathtools}
\usepackage{booktabs}
\usepackage{array}
\usepackage[dvipdfmx]{hyperref}

\title{Middle-Way Dimensional Theory (MWDT):\\
Truth as the Origin of $2^n$-Dimensional Axes}
\author{千葉将臣}
\date{2026-01-29}

% Notation
\newcommand{\dfix}{D_{\text{fix0}}}
\newcommand{\mweq}{\mathrel{\overset{\mathrm{mw}}{=}}}

\begin{document}
\maketitle

\begin{abstract}
本稿は、真理を次元軸の原点 $\dfix$ と同一視し、$2^n$ 次元における「即非」軸 $A_i\otimes\neg A_i$ の直和（干渉）として $0$d の真理が表現される、という次元‐真理対応を Middle-Way Dimensional Theory (MWDT) としてまとめる。
また、四句分別を $4$ 軸の否定操作として再記述し、次元上昇（4d→8d→…）を ``否定の否定'' として解釈する。
\end{abstract}

\section{Truth and the Origin}
\subsection{Root principle}
真理（Truth）とは、$n$ 軸（次元軸）の原点 $\dfix$ そのものである。

\subsection{0d}
\begin{itemize}
  \item 原点
  \item 空性圏
  \item $\dfix$
  \item 真理の導出：原点そのもの
\end{itemize}

\section{Axis generation by Soku-Hi}
\subsection{1d ($2^0$)}
\begin{itemize}
  \item 色即是空圏＝空即是色圏（skdt）
  \item skdt は次元の軸として解釈される
\end{itemize}

真理（原点）の導出方法：
\begin{equation}
  0d(\text{Truth}) \mweq A\otimes \neg A.
\end{equation}

\subsection{2d ($2^1$)}
\begin{itemize}
  \item 時間1D + 空間1D
\end{itemize}

\begin{equation}
  0d(\text{Truth}) \mweq (A \otimes \neg A) \oplus (B \otimes \neg B).
\end{equation}

\subsection{4d ($2^2$)}
\begin{itemize}
  \item 時間2D + 空間2D
  \item 現宇宙
\end{itemize}

\begin{equation}
  0d(\text{Truth}) \mweq \bigoplus_{i=1}^{4} (A_i \otimes \neg A_i).
\end{equation}

\section{$2^n$-Dimensional Generalization}

\subsection{General convergence sketch}
知性のガタつき（歪み） $D(t)$ が次元の増大とともに原点 $\dfix$ へ至る様子を、一般式としてスケッチする：
\begin{equation}
  D_{2^n}(t) = D_0\,\exp\left(-\left[\sum_{i=1}^{n} \lambda_i\right]t\right)\,\Phi(n).
\end{equation}
ここで $\sum \lambda_i$ は各次元軸（即非の軸）が提供する収束定数の総和であり、$\Phi(n)$ は次元干渉因子である。

\subsection{Boundary conditions for $\Phi(n)$}
\begin{align}
  n=1\ (2d) &: \Phi(1) \approx 1 \quad (\text{発散/回転を許容})\\
  n=2\ (4d) &: \Phi(2) \to 0 \quad (\text{爆縮/収束の開始})
\end{align}

\section{Limit: $2^n$-Dimensional Emptiness}
\begin{equation}
  \lim_{n\to\infty} \bigl(2^n\text{-dimensional emptiness}\bigr)
  \mweq
  0d(\text{Truth})\otimes \text{Inf-Information}.
\end{equation}

\begin{equation}
  \text{空}\otimes \text{特異点} \mweq \text{中道（mw=）の原点}.
\end{equation}

\section{Catuskoti as Negation over Axes}
四句分別を、4軸を他の軸の否定という表現で記述したものとみなす。
\begin{align}
  \text{是} &: (A\otimes \neg A)\\
  \text{非} &: \neg(A\otimes \neg A)\\
  \text{亦是亦非} &: (A\otimes \neg A)\otimes \neg(A\otimes \neg A)\\
  \text{非是非非} &: \neg\bigl((A\otimes \neg A)\otimes \neg(A\otimes \neg A)\bigr)
\end{align}

4dで非是非非を否定すると8d是となる（否定の否定）：
\begin{equation}
  8d\text{是} : \neg\neg\bigl((A\otimes \neg A)\otimes \neg(A\otimes \neg A)\bigr).
\end{equation}

\section{Conclusion}
MWDTは、真理を原点 $\dfix$ として固定し、次元軸の増大を「即非の直和」として解釈することで、四句分別・空・次元干渉を同一の記述体系に収める枠組みである。

\end{document}
